% !PRINTED 617
\weberpage{617}
%<<CONT
cyclic groups to the symmetric and the alternating group.

The permutation arising from $s$ consists of a single cycle with an odd number of terms $(0, 1, 2 \ldots n-1)$ and therefore belongs to the alternating group. The substitution $t$ leaves the index 0 unchanged and yields for the remaining figures again a single cycle. For by $t$, 1 passes into $g$, $g$ into $g^2$, $g^2$ into $g^3$, etc., and since the powers $1, g, g^2 \ldots g^{n-2}$, apart from their order, coincide with the figures $1, 2 \ldots n-1$, $t$ corresponds to the cycle $(1, g, g^2 \ldots g^{n-2})$, which consists of an even number of terms and consequently belongs to the permutations of the second kind, and is therefore not contained in the alternating group. On the other hand $t^2$ is again contained in it.

From this the result follows:

\gesperrt{IV. The full metacyclic group is not a divisor of the alternating group. The greatest common divisor of the two groups is the half-metacyclic group.}

Every transitive divisor of the metacyclic group is itself metacyclic.

The condition for the metacyclic equations can be given various forms, which can be derived from what has gone before.

The full linear group is contained in the symmetric group as a divisor of index $\nu = 1 . 2 . 3 \ldots (n-2)$. A function $y$ of the $n$ variables $x_0, x_1 \ldots x_{n-1}$ belonging to the full linear group, which we may call a \gesperrt{metacyclic function}, therefore satisfies a resolvent $F(y)=0$ of degree $\nu$, whose coefficients are symmetric functions of the $x$.

If now for $x$ one substitutes the roots of a metacyclic equation $f(x)=0$, then $y$ becomes rational.

If conversely the function $y$ becomes rational by the substitution of the roots of an irreducible equation $f(x)=0$ for the $x$, while $F'(y)$ remains different from zero, then $f(x)$ is metacyclic; for then the group of $f(x)=0$ is certainly a divisor of the full metacyclic group, and therefore, since $f(x)$ is irreducible, itself metacyclic. But it also suffices for the algebraic solubility of $f(x)=0$ if the resolvent $F(y)=0$ has any rational root at all, which
%CONT>>
